
\textbf{Figure 1} (\texttt{figures/fig1\textbackslash \_gate\textbackslash \_metrics.png}): Gate evaluation results per classification scheme. Left panel: coverage rate per scheme against the 70\% gate threshold. Right panel: cross-group edge count per scheme against the 30-edge threshold. Scheme 3 (FoS-primary) classifies all 33 resolved papers but yields only 9 cross-group edges at proxy corpus scale. Schemes 1 and 2 classify 4 papers each (12\%), producing zero cross-group edges.

\textbf{Figure 2} (\texttt{figures/fig2\textbackslash \_dropout\textbackslash \_by\textbackslash \_venue.png}): Unresolved papers by ID type. 16 of 49 papers fail S2AG resolution because they are ACM Digital Library publications without arXiv preprints. Dropout is non-systematic with respect to venue group, introducing no directional bias into the citation direction measurement.

\textbf{Figure 3} (\texttt{figures/fig3\textbackslash \_edge\textbackslash \_heatmaps.png}): Directed citation edge count heatmaps for all three classification schemes. Scheme 3 (right panel): $ML\_NLP\rightarrow ML\_NLP$ = 42, $ML\_NLP\rightarrow HCI$ = 5, $HCI\rightarrow ML\_NLP$ = 4, $HCI\rightarrow HCI$ = 0. Asymmetry ratio = 0.107. Schemes 1 and 2 (left and center panels): all cells = 0 (insufficient classification coverage).

\textbf{Figure 4} (\texttt{figures/fig4\textbackslash \_venue\textbackslash \_pie.png}): Corpus composition by venue group under Scheme 1 (venue-string) classification. Only 4 of 33 resolved papers are classified, illustrating the 12\% coverage limitation of venue-string approaches on S2AG data.

\textbf{Figure 5} (\texttt{figures/corpus\textbackslash \_overview.png}): Overview of the huashen218 bidirectional alignment corpus structure as accessed via the public GitHub reading list. Approximately 130 linked papers yield 49 extractable identifiers and 33 S2AG-resolved papers.
